% ============================================
% PÁGINA DE RESUMO
% ============================================
\thispagestyle{empty}

% Espaçamento 1,5 linhas
\onehalfspacing

% Recuo de parágrafo
\setlength{\parindent}{1.25cm}

\begin{center}
    \large\textbf{RESUMO}
\end{center}

\vspace{0.5cm}

\noindent A disciplina de compiladores é componente essencial dos cursos de Ciência da Computação, mas os algoritmos de análise sintática são notoriamente difíceis de ensinar e aprender devido à sua natureza abstrata e aos múltiplos passos intermediários. Este trabalho apresenta a VANSI (Visualização de Algoritmos de Análise Sintática), uma ferramenta educacional interativa e multiplataforma desenvolvida para auxiliar no ensino dos algoritmos de análise sintática LL(1), SLR e LR(1). Construída com o \textit{framework Svelte} e empacotada para \textit{desktop} e mobile por meio do \textit{Tauri} e do \textit{Capacitor}, a VANSI oferece execução passo a passo dos algoritmos com animações, autômatos interativos, árvores sintáticas, tabelas de parsing e pseudocódigo com \textit{breakpoints}. A ferramenta foi avaliada por meio de um estudo de métodos mistos com 23 estudantes da Universidade Federal do Ceará, Campus Quixadá, combinando \textit{logs} automatizados de interação com questionários estruturados. Os resultados indicam alta usabilidade (média 4,2/5) e eficácia pedagógica (média 4,4/5), com 91\% dos participantes expressando satisfação e 90\% manifestando intenção de reutilização. A VANSI preenche lacunas identificadas em ferramentas existentes ao oferecer maior interatividade e visualização detalhada dos processos internos dos algoritmos.

\vspace{0.5cm}

\noindent\textbf{Palavras-chave:} compiladores. algoritmos de análise sintática. visualização de algoritmos. tecnologia educacional. Svelte.
